\section{Evaluation}
\label{sec:evaluation}

The benchmark results reveal three substantial limitations in the current
implementation that bound sensitivity and memory use. We discuss each in turn,
along with the current workarounds.

% ── Duplex identification ─────────────────────────────────────────────────────
\subsection{Duplex identification failure}
\label{sec:eval:duplex}

Across all 24 samples, the duplex identification rate was below 0.025\%. The
design intent is for the endpoint fingerprint comparison to match forward and
reverse template reads from the same DNA fragment. In practice, the comparison
fails because the two templates are reverse complements: the XOR Hamming distance
between a forward-strand fingerprint and its reverse complement is approximately
32 bits (near-maximum for a 32-bit fingerprint), far above the identity threshold
of 4 bits used to define a match.

The consequence is that almost all molecules are classified as simplex. The strand
bias filter in the caller (Fisher's exact test on simplex forward and reverse
counts) still operates but has less power: duplex molecules provide the strongest
evidence against strand artefacts, and losing them pushes the confidence model
toward classifying borderline calls as low confidence.

The fix requires comparing the forward fingerprint against the \emph{reverse
complement} of the reverse fingerprint (or equivalently, canonicalising both
before comparison). This is a single-line change in the assembly stage and is
the highest-priority correctness issue to resolve before any production use.

% ── Sensitivity ceiling ───────────────────────────────────────────────────────
\subsection{Sensitivity ceiling at 200K reads}
\label{sec:eval:sensitivity}

The 200K read limit imposed for memory safety caps sensitivity at the sub-1\%
VAF regime. At 0.1\% VAF, a 200K-read sample contains on average 20 variant
read pairs targeting each 100\,bp window—after UMI grouping into molecules and
further compression by duplex collapsing, approximately 2–5 variant molecules
are expected per target. The current minimum alt molecule threshold is 2, so
borderline targets will be dropped. At 0.01\% VAF, fewer than one variant
molecule is expected in 200K reads: the assay simply cannot detect at this
depth.

The underlying bottleneck is the O($n^2$) UMI grouping step, which compares
all UMI pairs within a sample. For 1 million read pairs the peak RSS reaches
1.6\,GB; for 2 million it is estimated above 4\,GB. Switching to a hash-partition
strategy (group reads by UMI prefix, then compare within partitions) would reduce
this to O($n$) expected memory and allow millions of read pairs to be processed
without memory pressure.

% ── Graph coverage and false positives ───────────────────────────────────────
\subsection{Graph coverage and false positives}
\label{sec:eval:graph}

The pathfind stage identified variant paths for 258 of 375 targets (68.8\%) in
the 2\% VAF samples. The remaining 117 targets produced no alternative path.
Most of these are likely targets with very low coverage in the 200K-read subset,
or targets where the variant k-mer is present but the graph walk exhausts its
depth limit before finding a valid end anchor.

False positives (98--101 per sample at 2\% VAF) arise mainly from spurious
paths: the de Bruijn graph contains k-mer paths that differ from the reference
but arise from sequencing error clusters or from reads spanning nearby germline
variants. Raising the minimum alt molecule count from 2 to 3 would reduce false
positives at the cost of some true positive sensitivity; this trade-off is
sample-depth-dependent and could be tuned per application.

% ── Overall assessment ────────────────────────────────────────────────────────
\subsection{Overall assessment}

At 2\% VAF with 200K reads, kam detects 40--48\% of truth variants with
precision of 0.60. This is a meaningful baseline for an alignment-free approach
on a 375-target panel without any alignment step. The dominant loss factors
are duplex orientation (causing near-total loss of duplex signal), shallow
depth (200K reads is insufficient for sub-1\% VAF), and k-mer graph artefacts.

Fixing the duplex orientation bug is expected to substantially improve both
sensitivity (more molecules passing the duplex threshold) and precision
(fewer strand-biased artefacts passing the strand-bias filter). Increasing
input read depth to 1 million or more reads, achievable once the O($n^2$) memory
issue is resolved, would extend detection to 0.1--0.5\% VAF. These are the two
highest-priority improvements before a production evaluation.
